% =========================================================================
%  Figura (ejemplo modelo): tres esferas concéntricas a, 2a, 3a con cargas
%  q, 2q, 3q. Potencial (en unidades de kq/a) antes y después de conectar la
%  interior a tierra:
%   antes:   3 (r<=1), 2 + 1/r (1..2), 1 + 3/r (2..3), 6/r (r>=3)
%   después: 0 (r<=1), 2 - 2/r (1..2), 1 (2..3), 3/r (r>=3)
% =========================================================================
\begin{figure}[t]
\centering
\begin{tikzpicture}[font=\small]
  % ---------------------------------------------------------------- (a)
  \begin{scope}[x=0.62cm, y=0.62cm]
    \draw[azulprof, line width=0.8pt] (0,0) circle (3);
    \draw[azulprof, line width=0.8pt] (0,0) circle (2);
    \shade[ball color=azulvelo] (0,0) circle (1);
    \draw[azulprof, line width=0.6pt] (0,0) circle (1);
    \node[font=\scriptsize] at (0,0.25) {$q$};
    \node[font=\scriptsize] at (0,-0.3) {$a$};
    \node[font=\scriptsize, fill=white, inner sep=0.5pt] at (90:2) {$2q$};
    \node[font=\scriptsize, fill=white, inner sep=0.5pt] at (90:3) {$3q$};
    \draw[azulprof, line width=0.5pt] (1,0) -- (3.6,0) -- (3.6,-0.5);
    \pic at (3.6,-0.5) {tierra};
    \node[fondo, font=\scriptsize, anchor=south, text=grispizarra] at (2.3,0.05) {hilo};
    \node at (0,-4.1) {(a)};
  \end{scope}
  % ---------------------------------------------------------------- (b)
  \begin{scope}[xshift=4.4cm, yshift=-2.25cm]
    \begin{axis}[tema, width=8.4cm, height=5.0cm, xmin=0, xmax=5, ymin=0, ymax=3.3,
        xlabel={$r/a$}, ylabel={$\phi\,/\,(\kc q/a)$}, xtick={0,1,2,3,4,5}, ytick={0,1,2,3},
        legend style={at={(0.99,0.97)}, anchor=north east}]
      \addplot[azulprof, line width=1.2pt] coordinates {(0,3) (1,3)};
      \addlegendentry{antes}
      \addplot[azulprof, line width=1.2pt, domain=1:2, samples=30, forget plot] {2+1/x};
      \addplot[azulprof, line width=1.2pt, domain=2:3, samples=30, forget plot] {1+3/x};
      \addplot[azulprof, line width=1.2pt, domain=3:5, samples=30, forget plot] {6/x};
      \addplot[azulcielo, line width=1.2pt, densely dashed] coordinates {(0,0) (1,0)};
      \addlegendentry{interior a tierra}
      \addplot[azulcielo, line width=1.2pt, densely dashed, domain=1:2, samples=30, forget plot] {2-2/x};
      \addplot[azulcielo, line width=1.2pt, densely dashed, forget plot] coordinates {(2,1) (3,1)};
      \addplot[azulcielo, line width=1.2pt, densely dashed, domain=3:5, samples=30, forget plot] {3/x};
      \foreach \x in {1,2,3}{\edef\tmp{\noexpand\draw[punteado] (axis cs:\x,0) -- (axis cs:\x,3.3);}\tmp}
      \draw[{Stealth[length=3pt]}-{Stealth[length=3pt]}, cargapos] (axis cs:3.12,1) -- (axis cs:3.12,2)
        node[midway, right, font=\scriptsize] {$\Delta\phi_c$};
    \end{axis}
    \node at (3.6,-1.4) {(b)};
  \end{scope}
\end{tikzpicture}
\caption{Ejemplo modelo. (a) Tres esferas concéntricas de radios $a$, $2a$ y $3a$; la interior se conecta a tierra. (b) Potencial antes y después de la conexión:
la esfera interior pasa de $q$ a $-2q$, el potencial de la externa baja de $2\kc q/a$ a $\kc q/a$ y, entre $2a$ y $3a$, el campo se anula porque la carga
encerrada es $-2q+2q=0$.}
\label{fig:sem-modelo}
\end{figure}
